\documentclass[a4paper,11pt,fleqn]{article}%fleqn pour aligner les equations à gauche
\usepackage{../../Styles/Cours}

\newcommand{\chnum}{Chapitre 7}
\newcommand{\chtitle}{Taux d'avancement d'une réaction}
\newcommand{\dtype}{TP}

\begin{document}
\maketitle

L'acide benzoïque est un solide blanc parfois utilisé comme conservateur alimentaire sous le code E210. Il s'agit d'un acide faible.\\
\textbf{Comment déterminer l'avancement final d'une réaction acide-base ?}\\

\noindent
\begin{tabular}{|m{9cm} m{9cm}|}
\hline
\textbf{Document 1 : Taux d'avancement d'une réaction} & \textbf{Document 2 : Programme Python}\\
Le taux d'avancement d'une réaction $\tau$ est défini comme le rapport de l'avancement final $x_f$ sur l'avancement maximal $x_{max}$ de cette réaction :
\begin{eqnarray}
\tau = \dfrac{x_f}{x_{max}} \nonumber
\end{eqnarray}
Le taux d'avancement d'une réaction dépend des conditions initiales et de la constante d'équilibre de la réaction, c'est-à-dire la constante d'acidité du couple acide-base, notée $K_A$, lors d'une réaction acide-base.

\vspace{1em}
\textbf{Document 3 : Acide benzoïque}

Formule semi-développée de l'acide benzoïque : \begin{center} \chemfig[atom sep = 2em]{*6(HC-CH=CH-C(-[0]C(=[7]O)-[1]OH)=CH-CH=)} \end{center} 
 & \begin{lstlisting}[language=Python]
import numpy as np
print("L'equation a resoudre est c_A*tau^2 + K_A*tau - K_A = 0")
c_A = float(input('c_A ='))
pK_A = float(input('pK_A ='))
K_A = 10**(-pK_A)
Delta = K_A**2 - 4*c_A*(-K_A)
if (Delta > 0) :
  x_1 = (-K_A - np.sqrt(Delta)) / (2*c_A)
  x_2 = (-K_A + np.sqrt(Delta)) / (2*c_A)
  print ("Deux solutions : tau_1 =" + str(x_1) + " et tau_2 = " + str(x_2))
elif (Delta == 0) :
  x = - K_A/(2*c_A)
  print ("Une solution double : tau = ", + str(x))
else : 
  print ("Aucune solution")
print ("fin")
\end{lstlisting}\\
\hline
\end{tabular}

\noindent
\fbox{
\begin{minipage}{.975\textwidth}
\textbf{Données :}
\begin{itemize**}
	\item Solubilité de l'acide benzoïque dans l'eau à \num{20} \unit{\celsius} : $s$ = \SI{2,9}{\gram \per \liter}
	\item Couple acide-base : \ce{C6H5COOH(aq)}/\ce{C6H5COO-(aq)}
	\item $pK_A$ du couple considéré : $pK_A$ = 4,2
	\item Soit l'équation du seconde degré : $a\cdot x^2 + b \cdot x + c =0$ a pour discriminant $\Delta = b^2 - 4 \cdot a \cdot c$ :
	\begin{itemize}
		\item si $\Delta < 0$, l'équation n'a pas de solution réelle
		\item si $\Delta =0$, l'équation a une seule solution : $x=-\dfrac{b}{2a}$
		\item si $\Delta >0$, l'équation a deux solutions  : $x_1=\dfrac{-b-\sqrt{\Delta}}{2a}$ et $x_2=\dfrac{-b+\sqrt{\Delta}}{2a}$
	\end{itemize}
\end{itemize**}
\end{minipage}}

\textbf{Questions :} \\
\begin{enumerate}
	\item Écrire l'équation de la réaction entre l'acide benzoïque et l'eau
	\item Déterminer la concentration en quantité de matière d'acide benzoïque apporté $c_A$ d'une solution saturée en acide benzoïque
	\item Déterminer l'avancement maximal $x_{max}$ de la réaction de l'acide benzoïque avec l'eau dans un volume $V$ = \num{40} \unit{\milli \liter} de solution saturée.
	\item Cette solution a un pH de 2,9. Déterminer l'avancement final $x_f$, puis le taux d'avancement $\tau_{max}$ final de la réaction.
	\item \begin{enumerate}
		\item Écrire l'expression du quotient de réaction $Q_r$
		\item Montrer qu'à l'équilibre l'expression de $K_A$ d'écrit : $K_A=\dfrac{c_A \cdot \tau^2}{c^0 \cdot (1-\tau)}$
		\item Calculer $\tau$ avec le programme Python et comparer le résultat obtenu avec la valeur trouvée à la question \textbf{4}
	\end{enumerate}
\end{enumerate}
\textbf{Synthèse : } Expliquer en quoi la résolution de l'équation du second degré permet de calculer le pH d'une solution préparée par dissolution d'un acide faible dans l'eau.
\end{document}